% !PRINTED 645
\weberpage{645}
6. \gesperrt{Sind zwei ganze Funktionale mit einem dritten assoziiert, so sind sie auch untereinander assoziiert.}

7. \gesperrt{Die Norm $N(\alpha)$ ist durch $\alpha$ teilbar.}

Dies folgt aus der Gleichung (6), \S\,178: $\varphi(\alpha)=0$, deren letzter Koeffizient $\overline{A}_n=\pm N(\alpha)$ ist. Denn danach ist:
\[
\pm N(\alpha) = \alpha(\alpha^{n-1}+\overline{A}_1\alpha^{n-2}+\cdots+\overline{A}_{n-1}).
\]

8. \gesperrt{Ein ganzes Funktional, dessen absolute Norm eine Konstante ist, ist eine Einheit, und umgekehrt ist die absolute Norm einer Einheit $\varepsilon$ eine von Null verschiedene Konstante.}

Denn wenn $\varepsilon$ eine Einheit ist, so sind
\[
N(\varepsilon) \text{ und } N\left(\frac{1}{\varepsilon}\right) = \frac{1}{N(\varepsilon)}
\]
ganze Funktionale in $\overline{Z}$; folglich $N(\varepsilon)$ eine Einheit in $\overline{Z}$. Umgekehrt ist $N(\varepsilon)$ durch $\varepsilon$ teilbar, also, wenn $N(\varepsilon)$ eine Einheit in $\overline{Z}$ ist, $\varepsilon$ ein Teiler von 1, d.\,h. eine Einheit.

9. \gesperrt{Es gibt ganze rationale Funktionen von $z$, z.\,B. die absolute Norm von $\alpha$, die durch $\alpha$ teilbar sind. Ist $a$ eine durch $\alpha$ teilbare ganze rationale Funktion von $z$ von möglichst niedrigem Grade, so ist jede andere durch $\alpha$ teilbare ganze rationale Funktion von $z$ durch $a$ teilbar.}

\begin{center}
\S\,181. \gesperrt{Größter gemeinschaftlicher Teiler.}
\end{center}

Sind $\alpha, \beta, \ldots$ ganze Funktionale in $\overline{\Omega}$ und $x, y, \ldots$ Variable, die in $\alpha, \beta, \ldots$ nicht vorkommen, so ist
\weq{1}{\delta = \alpha x + \beta y + \cdots}
nach \S\,180, 2. ebenfalls ein ganzes Funktional, und zwar ist $\delta$ \gesperrt{teilbar durch jeden gemeinsamen Teiler von $\alpha, \beta, \ldots$} Sind $x_0, y_0, \ldots$ beliebige ganze Funktionen oder Funktionale in $\overline{Z}$, so ist, wie jetzt bewiesen werden soll,
\weq{2}{\delta_0 = \alpha x_0 + \beta y_0 + \cdots}
durch $\delta$ teilbar. Denn bezeichnen wir die absolute Norm von $\delta$ mit $D$, so ist
\weq{3}{N(\delta) = D E(x, y, \ldots) = D E,}
und $E(x, y, \ldots)$ ist eine Einheit in $\overline{Z}$, zugleich aber eine ganze homogene Funktion der Variablen $x, y, \ldots$
